\documentclass[../../../../SoftwareRequirementsSpecification.tex]{subfiles}
\begin{document}
% Richiede le macro definite in src/macros/useCaseMacros.tex
% (incluse nel preambolo di SoftwareRequirementsSpecification.tex)

\ucstart
\begin{xltabular}{\textwidth}{|l|X|}
  \hline
  \ucid{PT-04} \\ \hline
  \ucname{Configurazione Sessioni di Allenamento in una
    Scheda} \\ \hline
  \ucactor{PT} \\ \hline
  \ucdescription{Il PT inserisce o modifica le sessioni di
    allenamento presenti in una scheda.} \\ \hline
  \ucprecond{Il PT deve essere autenticato. Deve esistere almeno una
    scheda da modificare. La scheda non deve avere assegnazioni
    attive.} \\ \hline
  \ucpostcond{Il sistema configura le sessioni di allenamento
    della scheda.} \\ \hline
  \ucmainflow{
    \item Il PT apre la pagina di configurazione delle
      sessioni di allenamento.
    \item Il sistema mostra la lista di sessioni di allenamento presenti.
    \item Il PT preme su una delle sessioni di allenamento.
    \item Il sistema mostra il dettaglio della sessione di allenamento
      corrente.
    \item Il PT preme il pulsante "Aggiungi Esercizio".
    \item Il sistema mostra una lista di esercizi.
    \item Il PT sceglie un esercizio dalla lista.
    \item Il sistema mostra il form di inserimento dettagli esercizio.
    \item Il PT configura l'esercizio inserendo:
      \begin{itemize}
        \item Numero di serie/ripetizioni.
        \item Tempo di recupero tra una serie e l'altra (in secondi).
        \item Variante di esecuzione se presente (superserie, dropset,
          rest-pause etc...).
        \item Altre note esecutive.
      \end{itemize}
      Il carico non viene chiesto qui: è specifico dell'atleta e
      viene inserito al momento dell'assegnazione, nel caso d'uso
      PT-05 (vedi le Note).
    \item Il PT preme il pulsante "Aggiungi Esercizio".
    \item Il sistema aggiunge l'esercizio alla sessione.
    \item Il flusso riprende dal passo 4.
  } \\ \hline
  \ucaltflow{
    \textbf{5a. Modifica esercizio} \newline
    - Invece di aggiungere un nuovo esercizio, il PT preme su
    un esercizio già presente e lo modifica. \newline
    - Il flusso riprende dal passo 4. \newline
    \textbf{5b. Modifica sessione scelta} \newline
    - Il PT decide di modificare un'altra sessione di
    allenamento invece di quella corrente.\newline
    - Il flusso riprende dal passo 1.
  } \\ \hline
  \ucvariations{
    \textbf{9a. Selezione variante di esecuzione} \newline
    - Se il PT seleziona una variante, il sistema può chiedere dati
    aggiuntivi. Ad esempio: altri esercizi da eseguire in superserie,
    oppure il tempo di recupero tra una miniserie e l'altra nel
    rest-pause.
  } \\ \hline
  \ucnote{
    \textbf{Nota 1:} la scheda definisce \emph{quali} esercizi si
    svolgono e \emph{come} si svolgono (serie/ripetizioni, recupero,
    variante), ma non con quale peso. Il peso è specifico dell'atleta
    e viene inserito al momento dell'assegnazione, nel caso d'uso
    PT-05. È questo che permette di assegnare la stessa scheda a più
    atleti, come richiede lo statement.\newline
    \textbf{Nota 2:} la scheda non contiene alcuna regola di
    progressione del carico. Il carico può variare da una sessione
    all'altra perché il PT lo inserisce sessione per sessione in
    PT-05, non perché il sistema lo calcoli. Questo è coerente con lo
    statement, che ammette che il peso e le serie/ripetizioni variino
    da una sessione di allenamento all'altra senza indicare come.
  } \\ \hline
\end{xltabular}
\end{document}
